% Section 4, part 1 of 4; label prefix: tra. Part 2: frames, the affine slice, bivectors.
% Part 3: pointer to the u-adic convergence lemma. Part 4: dressings, edges, Lemmas E, T, F, the canonical gauge.
% No new macros. Local letters (declared where used): phi (a map of finite type), chi (sequences), pi_r, N, Theta.

\section{The vanishing tower and transport}
\label{sec:tra}

This section fixes the two coefficient settings in which the model is used, defines the vanishing tower $\Pi_i$, the slice $\sK$ and its frames, and proves the transport lemmas: for every symplectic dressing $P$ of the model that annihilates the vanishing period, in particular for the canonical gauge $\Pc$ of \cref{def:gau-can}, the edges lie in $\sKo\otimes\sKo$, the leaf point lies in $\sK$, its dilaton rescaling lies in the affine slice $-z+\sK$, and every graph term is a finite sum. Every infinite sum is justified $u$-adically. In this section and the next, $q$ denotes an element of a loop space, as in Givental's formalism; the parameter $q$ of the formal quintic does not occur.

\subsection{The vanishing tower and the slice}
\label{sec:tra-tower}

\begin{conv}[loop spaces, the two settings, forms, insertions]
\label{conv:tra-forms}
The state space $\bbH$ has the flat basis $H^a$, $a\in\bbZ/5$, with unit $1 = H^0$ and $H = H^1$, and the pairing $\eta(H^a,H^b) = 5\,\delta_{a+b\equiv 3}$ of \cref{def:set-fq}; as a matrix, $\eta$ is $5$ times the permutation matrix of $a\mapsto 3-a$, so $\eta^{-1} = \eta/25$. The charges and signs are those of \cref{def:set-model}, $\nu = (\nu_0,\dots,\nu_4) = (0,-1,3,2,1)$ and $\varepsilon = (1,-1,-1,1,1)$; put $w_b := \nu_b+1$. The \emph{weight} of $v^mz^kH^b$ is $5m+4k+w_b$ and the weight of a scalar $v^mz^k$ ($k\in\bbZ$) is $5m+4k$; $V_w$ denotes the $\bbQ$-span of the $v^mz^kH^b$ ($m,k\geq 0$) of weight $w$, which is the weight space of \cref{def:mod-weights}. In terms of the exponent weights of \cref{lem:set-modhom}, $4k+w_b = 5\wt(k,b)+5$.
\begin{description}
\item[Graded setting] Coefficients in $\bbQ[v]$. Each $V_w$ is finite-dimensional and consists of polynomials in $z$; put $\cH^{\mathrm{gr}}_+ := \bigoplus_wV_w = \bbH\otimes\bbQ[v][z]$. Givental's loop space $\bbH\otimes\bbQ[v]\lser{z^{-1}}$, with positive half $\cH^{\mathrm{gr}}_+$, is consistent in this setting, and every sum below is finite.
\item[$u$-adic setting] $v$ is a non-zero element of $u\bbQ\fps{u}$, written $v(u)$ when its dependence on $u$ is stressed; from \cref{def:tra-dressing} on, $v'(0)\neq 0$; for the formal quintic it is the series $v(u) = u\tau(u)^{-5/4}$ of \cref{def:set-tau}. The loop space and its halves are
\[
\cH^u := \bbH\otimes\bbQ\lser{z}\fps{u},\qquad \cH^u_+ := \bbH\otimes\bbQ\fps{u,z},\qquad \cH^u_- := \bbH\otimes z^{-1}\bbQ[z^{-1}]\fps{u}.
\]
Here $\bbQ\lser{z}\fps{u}$ is the power-series ring in $u$ over the field $\bbQ\lser{z}$: its elements are the series $\sum_{k\in\bbZ}c_kz^k$ with $c_k\in\bbQ\fps{u}$ and $c_k\to 0$ $u$-adically as $k\to-\infty$, and products are finite modulo every power of $u$. It contains the one-variable objects of this section: the leaf points and the entries of the dressings are power series in $z$, and the vanishing solution has infinitely many negative powers of $z$; the edges, in two variables, are treated in \cref{lem:tra-twoslot,lem:tra-E}. Every $f\in\cH^u$ splits uniquely, coefficientwise in $u$, as $f = f_{\geq 0}+f_{<0}$ with $f_{\geq 0}\in\cH^u_+$ and $f_{<0}\in\cH^u_-$; $f_{\geq k}$ is the part with $z$-powers $\geq k$. The form
\[
\Omega_{\mathrm s}(f,g) := \Res_z\,\eta(f(-z),g(z))\,\dd z,
\]
with the residue taken coefficientwise in $u$, is $\bbQ\fps{u}$-bilinear and antisymmetric on $\cH^u$ (the substitution $z\mapsto -z$ is an automorphism of $\bbQ\lser{z}$ which changes the sign of the residue, and $\eta$ is symmetric), and $\cH^u_+$, $\cH^u_-$ are isotropic and complementary. The substitution $v\mapsto v(u)$ maps every $V_w$ into $\cH^u_+$, and every scalar series $\sum_{n\geq 1}c_nv^{(w+4n)/5}z^{-n}$ of scalar weight $w$ ($c_n\in\bbQ$) into $z^{-1}\bbQ[z^{-1}]\fps{u}$, since its $z^{-n}$-coefficient has $u$-order at least $(w+4n)/5$.
\end{description}
\emph{Insertions and components.} An insertion $\psi^kH^b$ is identified with the vector $z^kH^b$. For $q$ in $\cH^{\mathrm{gr}}_+$ or $\cH^u_+$ we write $q_{(k,b)} := \coef{z^k}q^b$ for its flat components, so $q = \sum_{k,b}q_{(k,b)}z^kH^b$. Solutions are written in lowered components: a vector $J$ has the components $Y_a := \eta(H^a,J)$, i.e.\ $Y = \eta J$, and for $n = \sum_an^aH^a$ we put $\langle Y,n\rangle := \sum_aY_an^a = \eta(J,n)$. Matrices $P$ (gauges and dressings) act on lowered components; the \emph{vector form} of $P$ is $P^{\mathrm{vec}} := \eta^{-1}P\eta$, the conjugate of $P$ by the permutation $a\mapsto 3-a$, which acts on vectors. Since $\eta^{-1} = \eta/25$, the condition $P(-z)^T\eta P(z) = \eta$ holds for $P$ if and only if it holds for $P^{\mathrm{vec}}$. The vanishing solution $\YC$ of \cref{def:gau-van} has the vector form $J^{\mathrm C} := \eta^{-1}\YC$; thus $\eta(J^{\mathrm C}(-z),q(z)) = \langle\YC(-z),q(z)\rangle$.

\emph{Summable families; maps of finite type.} A family $(x_l)_{l\in L}$ in $\cH^u_+$ is \emph{summable} if for every flat index $(k,b)$ and every $N$ only finitely many $x_l$ have $\coef{z^k}x_l^b\notin u^N\bbQ\fps{u}$; its sum is taken coefficientwise and is finite modulo every power of $u$. The same definition is used for bivectors in $\bbH\otimes\bbH\otimes\bbQ\fps{\psi,\psi'}\fps{u}$, with the indices $(k,l,a,b)$ of the coefficients of $\psi^k\psi'^lH^a\otimes H^b$. A family $(y_l)$ in $\bbQ\lser{u}$ is summable if all $y_l$ lie in one $u^{N_0}\bbQ\fps{u}$ and for every $N$ only finitely many lie outside $u^N\bbQ\fps{u}$. Every regrouping of a summable family (a partition into subfamilies) consists of summable families; the family of their sums is summable (for a fixed coefficient index and $N$ the finitely many members outside $u^N\bbQ\fps{u}$ lie in finitely many subfamilies, and every other sub-sum lies in $u^N\bbQ\fps{u}$); and the sum of their sums is the total sum, because coefficientwise and modulo $u^N$ this is a statement about finite sums; the sum of a summable family of elements of $\bbQ\fps{u}$ lies in $\bbQ\fps{u}$. A $\bbQ\fps{u}$-linear map $\phi\colon\cH^u_+\to\bbQ\lser{u}$ (a local letter) is \emph{of finite type} if $\phi(\cH^u_+)$ lies in some $u^{N_0}\bbQ\fps{u}$ and for every $N$ there are a finite set $F_N$ of flat indices and an integer $M_N$ such that $\phi(x)\in u^N\bbQ\fps{u}$ whenever $\coef{z^k}x^b\in u^{M_N}\bbQ\fps{u}$ for all $(k,b)\in F_N$, the other coefficients of $x$ being arbitrary; the same definition is used for maps on bivectors.
\end{conv}

\begin{lem}[the summation lemma]\genusstatus{AI-proposed, unchecked}
\label{lem:tra-sum}
If $\phi$ is of finite type and $(x_l)_{l\in L}$ is a summable family in $\cH^u_+$, then $(\phi(x_l))_{l\in L}$ is summable and $\phi(\sum_lx_l) = \sum_l\phi(x_l)$. The same holds for maps of finite type on bivectors and summable families of bivectors.
\end{lem}

\begin{proof}
Fix $N$, with $F_N$, $M_N$ as in \cref{conv:tra-forms}. By summability only a finite set $L_0\subset L$ of indices has a coefficient $\coef{z^k}x_l^b$, $(k,b)\in F_N$, outside $u^{M_N}\bbQ\fps{u}$; so $\phi(x_l)\in u^N\bbQ\fps{u}$ for $l\notin L_0$, and, as $\sum_{l\notin L_0}x_l$ has all its coefficients with index in $F_N$ in $u^{M_N}\bbQ\fps{u}$, linearity gives $\phi(\sum_lx_l)-\sum_{l\in L_0}\phi(x_l)\in u^N\bbQ\fps{u}$. Since all $\phi(x_l)$ lie in $u^{N_0}\bbQ\fps{u}$, $(\phi(x_l))$ is summable, and $\phi(\sum_lx_l)\equiv\sum_l\phi(x_l)$ modulo every $u^N$. Bivectors are treated in the same way.
\end{proof}

The maps of finite type used below are $x\mapsto\coef{z^k}x^b$, the $\Pi_i$, a model correlator in one insertion (\cref{lem:tra-conv}), $\Omega_{\mathrm s}(f,\cdot)$ for $f\in\cH^u_-$, and, coefficientwise in the other slot, $\Pi_i\otimes 1$ (\cref{lem:tra-twoslot}) and the contraction of \cref{lem:tra-E}.

\begin{defn}[the vanishing tower]
\label{def:tra-tower}
Let $\YC$ be the vanishing solution of the model (\cref{def:gau-van}), in lowered components. For $q$ in $\cH^{\mathrm{gr}}_+$ or in $\cH^u_+$ and $i\geq 0$ put
\begin{equation}
\label{eq:tra-tower}
\Pi_i(q) := -\coef{z^i}\,\langle\YC(-z),q(z)\rangle .
\end{equation}
$\Pi_0$ is the \emph{vanishing covector}. In the $u$-adic setting $\YC$ is taken at $v = v(u)$, where it lies in $\bbQ[z^{-1}]\fps{u}^5$ (\cref{lem:gau-vanshape}), and the product is taken in $\bbQ\lser{z}\fps{u}$ (\cref{conv:tra-forms}).
\end{defn}

\emph{The shift rule and the values.} The series $\YC$ has no positive power of $z$. Since $\coef{z^i}(z^k\YC_b(-z)) = (-1)^{k-i}\coef{z^{-(k-i)}}\YC_b$,
\begin{equation}
\label{eq:tra-shift}
\Pi_i(z^kH^b) = \Pi_0(z^{k-i}H^b)\ \ (k\geq i),\qquad \Pi_i(z^kH^b) = 0\ \ (k<i),\qquad \Pi_0(z^mH^b) = (-1)^{m+1}\coef{z^{-m}}\YC_b .
\end{equation}
By the closed form of \cref{lem:gau-vanshape}, $\coef{z^{-n}}\YC_b = \varepsilon_b(-1)^nv^{n-r}/(4^r\,r!\,(n-r)!)$ with $r = \lfloor n/5\rfloor$ if $n\equiv\nu_b+1\pmod 5$, and $\coef{z^{-n}}\YC_b = 0$ otherwise. As $\nu_b+1\equiv 1-b\pmod 5$, this gives
\begin{equation}
\label{eq:tra-values}
\Pi_0(z^mH^b) = -\varepsilon_b\,\frac{v^{m-r}}{4^r\,r!\,(m-r)!}\quad\bigl(r = \lfloor m/5\rfloor\bigr)\ \ \text{if } b+m\equiv 1\pmod 5,\qquad \Pi_0(z^mH^b) = 0\ \ \text{otherwise}.
\end{equation}
In particular:
\begin{enumerate}
\item $\Pi_0(z^mH) = \pi_rv^{4r}$ if $m = 5r$, and $\Pi_0(z^mH) = 0$ if $5\nmid m$, where $\pi_r := 1/(4^r\,r!\,(4r)!)$, $\pi_0 = 1$;
\item $\Pi_0(H^a) = 0$ for $a\neq 1$;
\item the $v$-exponent of a non-zero $\Pi_0(z^mH^b)$ is $m-\lfloor m/5\rfloor\geq 4m/5$; so in the $u$-adic setting $\ord_u\Pi_0(z^mH^b)\geq 4m/5$.
\end{enumerate}
For example $\Pi_0(H) = 1$, $\Pi_0(z1) = -v$ and $\Pi_0(z^6\,1) = -v^5/480$.

\emph{Homogeneity and finite type.} The coefficient of $z^{-n}$ in $\YC_b$ is a multiple of $v^{n-\lfloor n/5\rfloor}$, of scalar weight $5(n-\lfloor n/5\rfloor)-4n = w_b$ when $n\equiv\nu_b+1\pmod 5$ (as $0\leq\nu_b+1\leq 4$); so $\YC_b$ is homogeneous of scalar weight $w_b$, and in the graded setting $\Pi_i$ maps $V_w$ into $\bbQ v^{(w-4i)/5}$ (zero unless $(w-4i)/5$ is a non-negative integer). In the $u$-adic setting $\Pi_i(q) = \sum_{k\geq i,\,b}q_{(k,b)}\,\Pi_0(z^{k-i}H^b)$ converges by (iii), and modulo $u^N$ it depends only on the $q_{(k,b)}$ with $4(k-i)/5<N$, taken modulo $u^N$: $\Pi_i$ maps $\cH^u_+$ to $\bbQ\fps{u}$ and is of finite type.

\emph{The tower as symplectic pairings.} For $k\geq 1$, $z^{-k}J^{\mathrm C}\in\cH^u_-$ by \cref{lem:gau-vanshape}, and for $q\in\cH^u_+$
\begin{equation}
\label{eq:tra-pairing}
\Omega_{\mathrm s}(z^{-k}J^{\mathrm C},q) = (-1)^{k+1}\,\Pi_{k-1}(q),
\end{equation}
because $(z^{-k}J^{\mathrm C})(-z) = (-1)^kz^{-k}J^{\mathrm C}(-z)$ and $\eta(J^{\mathrm C}(-z),q(z)) = \langle\YC(-z),q(z)\rangle$, so the left side is $(-1)^k\coef{z^{k-1}}\langle\YC(-z),q(z)\rangle$.

\begin{lem}[the operator $\Theta$]\genusstatus{AI-proposed, unchecked}
\label{lem:tra-theta}
Let $A$ be $\bbQ[v]$ (graded setting) or $\bbQ\fps{u}$ with $v = v(u)$ ($u$-adic setting). For a sequence $\chi = (\chi_j)_{j\geq 0}$ in $A$, finitely supported in the graded setting, put $\chi(z) := \sum_j\chi_jz^j$ and $\Theta\chi := (\Pi_i(\chi(z)H))_{i\geq 0}$. Then $\Theta = 1+N$, where
\[
(N\chi)_i := \sum_{r\geq 1}\pi_r\,v^{4r}\,\chi_{i+5r},
\]
and $\Theta$ is bijective on the set of these sequences, with inverse $\sum_{r\geq 0}(-N)^r$. The operator $N$, and hence $\Theta^{-1}$, maps sequences vanishing outside $[0,L)$ to sequences vanishing outside $[0,L)$; in the graded setting both map sequences with $\chi_j\in\bbQ v^{(w-4j)/5}$ for all $j$ to sequences of the same kind.
\end{lem}

\begin{proof}
By the shift rule \eqref{eq:tra-shift} and the value (i) above, $\Pi_i(\chi(z)H) =\sum_{j\geq i}\chi_j\Pi_0(z^{j-i}H) = \sum_{r\geq 0}\pi_rv^{4r}\chi_{i+5r}$, which is $\chi_i+(N\chi)_i$ (a convergent sum in the $u$-adic setting, as $\ord_uv^{4r}\geq 4r$). In both settings $(1+N)\sum_{r\leq R}(-N)^r\chi = \chi-(-N)^{R+1}\chi = \sum_{r\leq R}(-N)^r(1+N)\chi$. Graded setting: if $\chi$ vanishes outside $[0,L)$, then $N\chi$ vanishes outside $[0,L-5)$, so $N^r\chi = 0$ for $5r\geq L$ and the sums are finite. $u$-adic setting: $\ord_u(N^r\chi)_i\geq 4r+\min_j\ord_u\chi_j$ uniformly in $i$, so $((-N)^r\chi)_{r\geq 0}$ is summable coefficientwise, $N$ commutes with its sum (the double family $(\pi_kv^{4k}((-N)^r\chi)_{i+5k})_{k\geq 1,\,r\geq 0}$ is summable at each $i$), and $(-N)^{R+1}\chi\to 0$. So $\sum_r(-N)^r$ is a two-sided inverse of $\Theta$. The statement on supports holds for $N$ by its definition, and in the graded setting $v^{4r}\chi_{i+5r}$ has scalar weight $20r+(w-4i-20r) = w-4i$ if $\chi_j$ has weight $w-4j$ for all $j$; both properties pass to $\Theta^{-1}$.
\end{proof}

\begin{defn}[the slice $\sK$ and $\sKo$]
\label{def:tra-slice}
In either setting, with $A$ as in \cref{lem:tra-theta} and $\cH^A_+ := \cH^{\mathrm{gr}}_+$, resp.\ $\cH^A_+ := \cH^u_+$,
\begin{equation}
\label{eq:tra-slice}
\sK := \{q\in\cH^A_+ : \Pi_i(q) = 0\ \text{for all } i\geq 1\},\qquad \sKo := \sK\cap\ker\Pi_0 .
\end{equation}
Equivalently, $\sKo$ is the set of $n\in\cH^A_+$ with $\langle\YC(-z),n(z)\rangle$ in $z^{-1}\bbQ[z^{-1}]\fps{u}$ in the $u$-adic setting, and in $z^{-1}\bbQ[v]\fps{z^{-1}}$ in the graded setting. The \emph{affine slice} is $-z+\sK := -z1+\sK$; its description is \cref{cor:tra-slice}. When the setting must be named we write $\sK^A$, $\sKo^A$; the \emph{homogeneous} elements of $\sK$ are the elements of the $\sK\cap V_w$.
\end{defn}

% Section 4, part 2 of 4: frames, the affine slice, coefficients of frames, frame expansion of bivectors.
% Local letters: chi^q (the H-part of q), pr_{K_0}, f_j^{(k,a)} (the coefficients of frame(k,a)), n(c).

\begin{lem}[frames]\genusstatus{AI-proposed, unchecked}
\label{lem:tra-frames}
Let $A$ be as in \cref{lem:tra-theta}, and $A\{z\} := \bbQ[v][z]$ in the graded, $A\{z\} := \bbQ\fps{u,z}$ in the $u$-adic setting.
\begin{description}
\item[(a)] $\cH^A_+ = \sKo^A\oplus A\{z\}H$. In the graded setting $V_w = (\sKo\cap V_w)\oplus(A\{z\}H\cap V_w)$ for every $w$.
\item[(b)] For every $(k,a)$ with $a\neq 1$ there is exactly one element of $\sKo^A$ of the form $z^kH^a+\chi(z)H$ with $\chi(z)\in A\{z\}$; it is
\[
\fr(k,a) = z^kH^a+\sum_{j<k}f_jz^jH,\qquad f_j = c_{k,a,j}\,v^{(4(k-j)+\nu_a+1)/5},\quad c_{k,a,j}\in\bbQ,
\]
where $c_{k,a,j} = 0$ unless the exponent is an integer, which is then $\geq 1$. No element of $\sKo^A$ has the form $z^kH+\sum_{j<k}\chi_jz^jH$.
\item[(c)] (Frame coordinates.) The map $\sKo^A\to\prod_{k\geq 0,\,a\neq 1}A$, $n\mapsto(n_{(k,a)})$, is injective. In the $u$-adic setting it is bijective; in the graded setting its image is the set of finitely supported families in $\bbQ[v]$, and on $\sKo\cap V_w$ the $(k,a)$-member is a rational multiple of $v^{(w-4k-\nu_a-1)/5}$. The inverse sends $(c_{(k,a)})$ to the sum of the summable family $(c_{(k,a)}\fr(k,a))$. In particular the projection $\mathrm{pr}_{\sKo}$ of $\cH^A_+$ onto $\sKo^A$ along $A\{z\}H$ is $\mathrm{pr}_{\sKo}(q) = \sum_{k,\,a\neq 1}q_{(k,a)}\fr(k,a)$.
\item[(d)] The frames of the $u$-adic setting are the frames of the graded setting at $v = v(u)$.
\end{description}
Moreover $\fr(k,a) = z^kH^a$ if and only if $k<\nu_a+1$; these ten \emph{trivial frames} are those with $(k,a) = (0,0)$, $(0,2)$, $(1,2)$, $(2,2)$, $(3,2)$, $(0,3)$, $(1,3)$, $(2,3)$, $(0,4)$, $(1,4)$. For example
\begin{gather*}
\fr(1,0) = z1+vH =: \se,\qquad \fr(2,4) = z^2H^4+\tfrac{v^2}{2}H,\qquad \fr(3,3) = z^3H^3+\tfrac{v^3}{6}H,\\
\fr(4,2) = z^4H^2-\tfrac{v^4}{24}H,\qquad \fr(6,0) = z^6\,1+vz^5H-\tfrac{v^5}{120}H ;
\end{gather*}
$\se$ is the Euler--dilaton vector.
\end{lem}

\begin{proof}
\emph{The decomposition (a).} For $q\in\cH^A_+$ put $\chi^q := \Theta^{-1}((\Pi_i(q))_{i\geq 0})$ (\cref{lem:tra-theta}) and $\mathrm{pr}_{\sKo}(q) := q-\chi^q(z)H$. By the definition of $\Theta$, $\Pi_i(\chi^q(z)H) = (\Theta\chi^q)_i = \Pi_i(q)$, so $\mathrm{pr}_{\sKo}(q)\in\sKo^A$ and $q = \mathrm{pr}_{\sKo}(q)+\chi^q(z)H$; in the graded setting $(\Pi_i(q))_i$ is finitely supported and $\chi^q(z)$ is a polynomial. If $\chi(z)H\in\sKo^A$, then $\Theta\chi = 0$, hence $\chi = 0$ by the injectivity of $\Theta$: the sum is direct. For $q\in V_w$, $(\Pi_i(q))_i$ vanishes beyond the $z$-degree of $q$ and $\Pi_i(q)\in\bbQ v^{(w-4i)/5}$ (the paragraph on homogeneity after \cref{def:tra-tower}); by \cref{lem:tra-theta} $\chi^q_j\in\bbQ v^{(w-4j)/5}$, so $\chi^q(z)H$ and $\mathrm{pr}_{\sKo}(q)$ lie in $V_w$. This proves (a).

\emph{Frames: (b) and (d).} For $a\neq 1$ put $\fr(k,a) := \mathrm{pr}_{\sKo}(z^kH^a)$. The sequence $(\Pi_i(z^kH^a))_i$, which is $(\Pi_0(z^{k-i}H^a))_i$, vanishes for $i>k$ by the shift rule \eqref{eq:tra-shift} and at $i = k$ by the value (ii) after \cref{def:tra-tower}; so it vanishes outside $[0,k)$, hence so does $\chi^{z^kH^a}$ (\cref{lem:tra-theta}), and $\fr(k,a) = z^kH^a+\sum_{j<k}f_jz^jH$ with $f_j := -\chi^{z^kH^a}_j$. Uniqueness: the difference of two such elements lies in $\sKo^A\cap A\{z\}H = 0$. An element $z^kH+\sum_{j<k}\chi_jz^jH\neq 0$ lies in $A\{z\}H$, hence not in $\sKo^A$, by (a). Exponents: in the graded setting $z^kH^a\in V_w$ with $w = 4k+\nu_a+1$, so by the proof of (a) $f_j\in\bbQ v^{(w-4j)/5} = \bbQ v^{(4(k-j)+\nu_a+1)/5}$, and $f_j = 0$ unless this exponent is a non-negative integer, i.e.\ unless $j\equiv k-\nu_a-1\pmod 5$; for $j\leq k-1$ and $a\neq 1$ (so $\nu_a\geq 0$) the exponent is at least $(4+1)/5 = 1$. For (d): the substitution $v\mapsto v(u)$ commutes with every $\Pi_i$ on polynomials in $z$ over $\bbQ[v]$, so a graded frame at $v(u)$ lies in the $u$-adic $\sKo$, has the form of (b), and is the $u$-adic frame by uniqueness.

By \cref{lem:tra-theta} the frame condition $\Pi_i(\fr(k,a)) = 0$ ($i\geq 0$) reads $f_i+(Nf)_i = -\Pi_0(z^{k-i}H^a)$, i.e.
\[
f_i = -\Pi_0(z^{k-i}H^a)-\sum_{r\geq 1}\pi_rv^{4r}f_{i+5r}\qquad(0\leq i<k),
\]
from which the examples follow with \eqref{eq:tra-values}; e.g.\ for $\fr(6,0)$, $f_5 = -\Pi_0(z1) = v$ and $f_0 = -\Pi_0(z^6\,1)-\pi_1v^4f_5 = -v^5/120$. Trivial frames: if $k\geq\nu_a+1$, the index $j_0 := k-\nu_a-1$ satisfies $0\leq j_0<k$ and $j_0+5r>k$ for $r\geq 1$ (as $\nu_a\leq 3$), so the recursion gives $f_{j_0} = -\Pi_0(z^{\nu_a+1}H^a) = \varepsilon_av^{\nu_a+1}/(\nu_a+1)!\neq 0$; if $k<\nu_a+1$, no $j<k$ satisfies $j\equiv k-\nu_a-1\pmod 5$, since $1\leq k-j\leq\nu_a$, so all $f_j$ vanish. Counting the $k<\nu_a+1$ for $a = 0,2,3,4$ ($\nu_a = 0,3,2,1$) gives the ten trivial frames.

\emph{Frame coordinates: (c).} Let $c = (c_{(k,a)})_{k\geq 0,\,a\neq 1}$ be a family in $A$, finitely supported in the graded setting, and write $f^{(k,a)}_j$ for the coefficient $f_j$ of $\fr(k,a)$. The coefficient $\coef{z^j}(\cdot)^b$ of $c_{(k,a)}\fr(k,a)$ is, for $b\neq 1$, $c_{(j,b)}$ if $(k,a) = (j,b)$ and $0$ otherwise; for $b = 1$ it is $c_{(k,a)}f^{(k,a)}_j$, of $u$-order at least $(4(k-j)+\nu_a+1)/5\geq(4(k-j)+1)/5$ by (b) and $\ord_uv\geq 1$. So the family $(c_{(k,a)}\fr(k,a))$ is summable; let $n(c)$ be its sum. As each $\Pi_i$ is of finite type, \cref{lem:tra-sum} gives $\Pi_i(n(c)) = 0$: $n(c)\in\sKo^A$, with $n(c)_{(k,a)} = c_{(k,a)}$ for $a\neq 1$. Conversely, for $n\in\sKo^A$ and $c_{(k,a)} := n_{(k,a)}$, $n-n(c)$ has no component with $a\neq 1$, so it lies in $\sKo^A\cap A\{z\}H = 0$ by (a). Hence $n\mapsto(n_{(k,a)})$ is bijective with inverse $c\mapsto n(c)$. For any $q$, $\mathrm{pr}_{\sKo}(q)$ and $\sum_{k,\,a\neq 1}q_{(k,a)}\fr(k,a)$ are elements of $\sKo^A$ with the components of $q$ for $a\neq 1$ ($\mathrm{pr}_{\sKo}$ changes only $H$-components); so they are equal, and
\begin{equation}
\label{eq:tra-decomp}
q = \sum_{k,\,a\neq 1}q_{(k,a)}\fr(k,a)+\chi^q(z)H,\qquad \chi^q_j = q_{(j,1)}-\sum_{k>j,\,a\neq 1}q_{(k,a)}f^{(k,a)}_j,
\end{equation}
the last sum being $u$-adically convergent (finite in the graded setting). In the graded setting an element of $\sKo$ is a polynomial in $z$, so its frame coordinates are finitely supported, and on $\sKo\cap V_w$ the coordinate $n_{(k,a)} = \coef{z^k}n^a$ is the coefficient of a monomial $v^mz^kH^a$ of weight $w$, a rational multiple of $v^{(w-4k-\nu_a-1)/5}$; as $\Pi_i$ maps $V_w$ to $\bbQ v^{(w-4i)/5}$, $\sKo\cap\cH^{\mathrm{gr}}_+ = \bigoplus_w(\sKo\cap V_w)$.
\end{proof}

\begin{cor}[the affine slice]\genusstatus{AI-proposed, unchecked}
\label{cor:tra-slice}
$\sK^A = \sKo^A\oplus AH$, the decomposition being $q = (q-\Pi_0(q)H)+\Pi_0(q)H$. Moreover $-z1\in\sK^A$, so the affine slice $-z+\sK$ is the set $\{q\in\cH^A_+ : \Pi_i(q) = 0,\ i\geq 1\}$, i.e.\ $\sK^A$ itself; its points are written uniquely as $q = -z1+n+\mu H$ with $n\in\sKo^A$ and $\mu\in A$ (a local letter), and
\[
\Pi_0(-z1+n+\mu H) = v+\mu .
\]
\end{cor}

\begin{proof}
By the shift rule $\Pi_i(H) = 0$ for $i\geq 1$, and $\Pi_0(H) = 1$; so $H\in\sK^A$, $AH\cap\sKo^A = 0$, and $q-\Pi_0(q)H\in\sKo^A$ for $q\in\sK^A$. By \eqref{eq:tra-values} and the shift rule, $\Pi_0(-z1) = v$, $\Pi_1(-z1) = -\Pi_0(1) = 0$ and $\Pi_i(-z1) = 0$ for $i\geq 2$. Hence $-z1\in\sK^A$, and for $q = -z1+y$, $y = n+\mu H\in\sK^A$, linearity gives $\Pi_0(q) = v+\mu$.
\end{proof}

\begin{lem}[coefficients of frames]\genusstatus{AI-proposed, unchecked}
\label{lem:tra-coeff}
Over $\bbQ\fps{u}$, every $n\in\sKo$ is the sum of the summable family $(n_{(k,a)}\fr(k,a))$, where $\fr(k,a)\in\sKo\cap V_{4k+\nu_a+1}$ is homogeneous over $\bbQ[v]$, evaluated at $v(u)$, and every non-zero coefficient $f_j$ of $\fr(k,a)$ has the $v$-exponent $\wt(k,a)-\wt(j,1) = (4(k-j)+\nu_a+1)/5\geq 1$; every $q\in\sK$ is such a sum plus $\Pi_0(q)H$. So every insertion used in \cref{thm:red-main} is a $\bbQ\fps{u}$-combination, in the sense of summable families, of homogeneous elements of $\sKo$ and of $H$.
\end{lem}

\begin{proof}
\Cref{lem:tra-frames}~(b)-(d) and \cref{cor:tra-slice}, with $\wt(k,a)-\wt(j,1) = ((4k+\nu_a-4)-(4j-1-4))/5$.
\end{proof}

\begin{lem}[frame expansion of bivectors]\genusstatus{AI-proposed, unchecked}
\label{lem:tra-twoslot}
For $E\in\bbH\otimes\bbH\otimes\bbQ\fps{\psi,\psi'}\fps{u}$ write $E^{ab}_{kl} := \coef{\psi^k\psi'^l}E^{ab}$ (first index: first slot), put $E_{\cdot,(l,b)} := \sum_{k,a}E^{ab}_{kl}\psi^kH^a$, an element of $\cH^u_+$ in the variable $\psi$, and
\[
(\Pi_i\otimes 1)E := \sum_{l,b}\Pi_i(E_{\cdot,(l,b)})\,\psi'^lH^b ;
\]
$(1\otimes\Pi_i)E$ is defined in the same way with the two slots exchanged. The following are equivalent: (i) $(\Pi_i\otimes 1)E = 0$ and $(1\otimes\Pi_i)E = 0$ for all $i\geq 0$; (ii) $E$ is the sum of the summable family $(E^{\alpha\beta}\,\fr_\alpha(\psi)\otimes\fr_\beta(\psi'))$ over $\alpha = (k,a)$, $\beta = (l,b)$ with $a,b\neq 1$, where $E^{\alpha\beta} := E^{ab}_{kl}$. The coefficients in (ii) are unique: they are the flat components of $E$ with $a,b\neq 1$.
\end{lem}

\begin{proof}
The family in (ii) is summable for any coefficients in $\bbQ\fps{u}$: the coefficient of $\psi^{k'}\psi'^{l'}H^{a'}\otimes H^{b'}$ of the $(\alpha,\beta)$-term is a product of a coefficient of $\fr_\alpha$ and one of $\fr_\beta$; for $a'\neq 1$ only $\alpha = (k',a')$ contributes, for $a' = 1$ the $u$-order is at least $(4(k-k')+1)/5$ (the exponent of \cref{lem:tra-frames}~(b), with $\nu_a\geq0$ and $\ord_uv\geq1$), and likewise in the second slot. The same computation shows that the coefficient of $\psi^k\psi'^lH^a\otimes H^b$, $a,b\neq 1$, of the sum is $E^{\alpha\beta}$: this is the uniqueness.
(ii)$\Rightarrow$(i): coefficientwise in the other slot, $\Pi_i\otimes 1$ and $1\otimes\Pi_i$ are of finite type and kill every term; apply \cref{lem:tra-sum}.
(i)$\Rightarrow$(ii): for each $(l,b)$, $E_{\cdot,(l,b)}\in\sKo$ by the first half of (i), so $E_{\cdot,(l,b)} = \sum_\alpha E^{ab}_{kl}\fr_\alpha(\psi)$ over $\alpha = (k,a)$, $a\neq 1$ (\cref{lem:tra-frames}~(c)); regrouping the double family $\bigl(E^{ab}_{kl}\,\fr_\alpha(\psi)\otimes\psi'^lH^b\bigr)$ over $(\alpha,(l,b))$, summable by the computation of the first paragraph in the first slot (the second slot is a flat monomial), gives the summable family $(\fr_\alpha\otimes G_\alpha)_\alpha$ (\cref{conv:tra-forms}) and $E = \sum_\alpha\fr_\alpha(\psi)\otimes G_\alpha(\psi')$ with $G_\alpha := \sum_{l,b}E^{ab}_{kl}\psi'^lH^b$. By \cref{lem:tra-sum}, $0 = (1\otimes\Pi_i)E = \sum_\alpha\Pi_i(G_\alpha)\fr_\alpha(\psi)$, an element of $\sKo$ with frame coordinates $(\Pi_i(G_\alpha))_\alpha$, which therefore vanish. So $G_\alpha\in\sKo$ and $G_\alpha = \sum_\beta E^{\alpha\beta}\fr_\beta(\psi')$; the resulting double family is summable, with sum $E$.
\end{proof}

``Edges in $\sKo\otimes\sKo$'' means (ii), a completed tensor product made precise by summable families.

% Section 4, part 3 of 4: the u-adic convergence lemma (stated and proved in an appendix).

All infinite sums of this section and of \cref{sec:red,sec:flh} are justified by \cref{lem:tra-conv}, stated and proved in \cref{sec:tra-convproof}: frames and their expansions inside correlators, edges and leaves are summable (parts (a), (b)); the vanishing series and its pairings are ring operations in $\bbQ\lser{z}\fps{u}$ (part (c)); leaves of $\psi$-degree $\geq 2$ are bounded by the dimension (part (d)); the vertex graph resums under the Model Lemma (part (e)); and the dilaton re-expansion of a vertex tensor converges when the dilaton factor is $1$ at $u = 0$, and in general not otherwise (part (f)).

% Section 4, part 4 of 4: dressings, edges and the leaf point; Lemmas E, T, F; the canonical gauge.
% Local letters: the ring cA = Q((psi))[[u,psi']], alpha, beta, C (in lem:tra-E); cN (numerator of the edge), xi, G.

\subsection{Dressings, edges and the leaf point}
\label{sec:tra-dress}

\begin{defn}[dressings]
\label{def:tra-dressing}
Let $v = v(u)\in u\bbQ\fps{u}$ with $v'(0)\neq 0$. A \emph{dressing of the model at $v(u)$} is a matrix series $P = \sum_{k\geq 0}z^kP_k(u)$ with $P_k\in\mathfrak{gl}_5(\bbQ\fps{u})$, graded, i.e.\ $P_k = \sum_bp_{k,b}E_{b+k,b}$ with the matrix units $E_{a,b}$ (indices modulo $5$), and with $P_0$ invertible; it acts on lowered components, $Y\mapsto PY$. We consider the two conditions
\begin{description}
\item[(Sp)] $P(-z)^T\eta P(z) = \eta$;
\item[(A)] $Y' := P\,\YC(v(u))$ has no positive power of $z$.
\end{description}
Put $T_P := -\coef{z^{-1}}Y'_0\in\bbQ\fps{u}$, the transported vanishing coordinate; $d := (P_0)_{33} = p_{0,3} = (P_0^{\mathrm{vec}})_{00}$, the dilaton factor, a unit of $\bbQ\fps{u}$; and $R := P_0^{-1}P = 1+O(z)$. Under (Sp), the equation at $z^0$ reads $P_0^T\eta P_0 = \eta$; hence $R$ satisfies (Sp) as well, and $p_{0,0}\,p_{0,3} = 1$. Examples: (i) the canonical gauge $\Pc$ (\cref{prop:tra-can}), and likewise $P\,P_{11}(0,z)^{-1}$ for a graded regular gauge $P$ from the pulled-back model to a self-dual one-pole system with equal products and equal residues (\cref{prop:gau-graded}~(iii), \cref{thm:gau-G}~(c) with \cref{lem:gau-vanq}~(c), \cref{prop:gau-symp}; the pulled-back model has the one-pole shape for every $v$ with $v'(0)\neq0$, by the computation of \cref{prop:gau-equalprod}~(i)); (ii) $DP$ with $P$ as in (i) and $D$ graded over $\bbQ\fps{u}$, $D_0$ invertible, with (Sp) and $DY'$ free of positive powers of $z$; these need not be gauges (\cref{rmk:mai-annih}).
\end{defn}

\begin{defn}[edges, the leaf point]
\label{def:tra-edge}
Let $R = 1+O(z)$ lie in $\mathfrak{gl}_5(\bbQ\fps{u,z})$, acting on vectors, with $R(-z)^T\eta R(z) = \eta$, and write $R^{-1}(\psi) := R(\psi)^{-1}$. The \emph{edge} of $R$ is the bivector (first index: first slot)
\[
E_R(\psi,\psi') := \frac{\eta^{-1}-R^{-1}(\psi)\,\eta^{-1}R^{-1}(\psi')^T}{\psi+\psi'} ;
\]
over $\bbQ\fps{u}$ the numerator is divisible by $\psi+\psi'$ (see below), so $E_R\in\bbH\otimes\bbH\otimes\bbQ\fps{u,\psi,\psi'}$, and $E_R(\psi,\psi')^T = E_R(\psi',\psi)$. For a dressing $P$ with (Sp) let $R^{\mathrm{vec}} = (P_0^{\mathrm{vec}})^{-1}P^{\mathrm{vec}}$ be the vector form of $R = P_0^{-1}P$, and put
\[
E_P := E_{R^{\mathrm{vec}}},\qquad q_T := (P^{\mathrm{vec}})^{-1}(-z1)\in\cH^u_+,\qquad \tilde q_T := d\,q_T ;
\]
$q_T$ is the \emph{leaf point}. The only graph sum attached to $P$ that enters the note is the \emph{flat sum of $P.\Omega^{\mathrm{mod}}_{v(u)}$} of \cref{def:red-flatsum}: the leg-free graph sum with the edge $E_P$ on every edge and, at a vertex of genus $g_v$ with $n_v$ half-edges, the $n_v$-th derivative of the model potential $\overline F_{g_v}$ at $\tilde q_T$ multiplied by $d^{2g_v-2+n_v}$. Its relation with the $R$-matrix action of \cite[Def.~2.2]{PPZ15} is stated after \eqref{eq:tra-leaf}.
\end{defn}

\emph{Divisibility.} Put $\cN(\psi,\psi') := \eta^{-1}-R^{-1}(\psi)\eta^{-1}R^{-1}(\psi')^T\in\mathfrak{gl}_5(\bbQ\fps{u}\fps{\psi,\psi'})$, the numerator of $E_R$; $R$ is invertible in $\mathfrak{gl}_5(\bbQ\fps{u,z})$ because $R(u,0) = 1$. Inverting $R(-z)^T\eta R(z) = \eta$ gives
\[
R^{-1}(\psi)\,\eta^{-1}R^{-1}(-\psi)^T = \eta^{-1},\qquad\text{i.e.}\qquad \cN(\psi,-\psi) = 0 .
\] In the variables $\psi$ and $\xi := \psi+\psi'$ (a local letter), $\cN$ becomes a power series vanishing at $\xi = 0$, hence divisible by $\xi$; so $\cN = (\psi+\psi')E_R$ with a unique $E_R\in\mathfrak{gl}_5(\bbQ\fps{u}\fps{\psi,\psi'})$. As $\eta^{-1}$ is symmetric, $\cN(\psi,\psi')^T = \cN(\psi',\psi)$, hence $E_R(\psi,\psi')^T = E_R(\psi',\psi)$. For a dressing, $E_P$ is also given by the same formula with $P^{\mathrm{vec}}$ in place of $R^{\mathrm{vec}}$, because $(P_0^{\mathrm{vec}})^T\eta P_0^{\mathrm{vec}} = \eta$.

\emph{The leaf point.} Let $P$ be a dressing with (Sp). By the grading $P_0$ is diagonal, so $P_0^{\mathrm{vec}}$ is diagonal with $(P_0^{\mathrm{vec}})_{00} = (P_0)_{33} = d$, and $(P_0^{\mathrm{vec}})^{-1}1 = d^{-1}1$. Since $P^{\mathrm{vec}} = P_0^{\mathrm{vec}}R^{\mathrm{vec}}$ and $R^{\mathrm{vec}} = 1+O(z)$,
\begin{equation}
\label{eq:tra-leaf}
q_T = -\frac zd\,(R^{\mathrm{vec}})^{-1}(z)1 = -\frac zd\,1+T_{\geq 2},\qquad T_{\geq 2} := d^{-1}\bigl(z1-z(R^{\mathrm{vec}})^{-1}(z)1\bigr)\in z^2\cH^u_+ ,
\end{equation}
so that $\tilde q_T = dq_T = -z(R^{\mathrm{vec}})^{-1}(z)1 = -z1+dT_{\geq 2}$ and $dT_{\geq 2} = \tilde q_T+z1 = \ttr_{R^{\mathrm{vec}}}(z)$ is the translation of \cite[Def.~2.13]{PPZ15} (the $T_b$ of \cite[Prop.~2.12]{PPZ15}); this computation uses neither (Sp) nor (A). Thus $E_P = E_{R^{\mathrm{vec}}}$ and $\tilde q_T = q_{R^{\mathrm{vec}}}$ are the edge and the leaf point (\cref{def:set-graphsum}) of the unit-preserving action $R^{\mathrm{vec}}.\Omega^{\mathrm{mod}}_{v(u)}$ of \cite[Defs.~2.2, 2.13]{PPZ15}; the vertex tensors of \cref{def:red-flatsum} are those of this action multiplied by $d^{2g_v-2+n_v}$, and for $P_0 = 1$ the flat sum is the leg-free graph sum of $R^{\mathrm{vec}}.\Omega^{\mathrm{mod}}_{v(u)}$.

\subsection{Lemmas E, T, F and the canonical gauge}
\label{sec:tra-lemmas}

\begin{lem}[Lemma E]\genusstatus{AI-proposed, unchecked}
\label{lem:tra-E}
Let $R = 1+O(z)$ lie in $\mathfrak{gl}_5(\bbQ\fps{u,z})$ with $R(-z)^T\eta R(z) = \eta$, and let $f\in\cH^u_-$, both in vector form. Then $Rf\in\cH^u$, the contraction of the covector $\Omega_{\mathrm s}(f,\cdot)$ into the first slot of $E_R$ (it pairs $\psi^me$ with $\eta(\coef{z^{-m-1}}f(-z),e)$) converges $u$-adically, and
\begin{equation}
\label{eq:tra-lemE}
(\Omega_{\mathrm s}(f,\cdot)\otimes 1)\,E_R(\psi') = -R(\psi')^{-1}\,(Rf)_{\geq 0}(\psi')\qquad\text{in }\bbH\otimes\bbQ\fps{u,\psi'} .
\end{equation}
Hence $E_R$ annihilates $\Omega_{\mathrm s}(f,\cdot)$ in the first slot, equivalently (by symmetry) in the second, if and only if $Rf\in\cH^u_-$. The computation takes place in the ring $\cA := \bbQ\lser{\psi}\fps{u,\psi'}$, the power-series ring in $u$ and $\psi'$ over the field $\bbQ\lser{\psi}$ ($\psi$ the variable of the first slot).
\end{lem}

\begin{proof}
\emph{The ring.} $\cA$ contains $\bbQ\fps{u,\psi,\psi'}$ (hence $E_R$, $R^{-1}(\psi)$, $R^{-1}(\psi')$ and $\cN$) and $\bbQ[\psi^{-1}]\fps{u}$, and $\psi+\psi' = \psi(1+\psi^{-1}\psi')$ is a unit of $\cA$, with inverse $\sum_{j\geq 0}(-\psi')^j\psi^{-j-1}$. The residue $\Res_\psi\colon\cA\to\bbQ\fps{u,\psi'}$, the coefficient of $\psi^{-1}$ taken coefficientwise in $(u,\psi')$, is $\bbQ\fps{u,\psi'}$-linear and commutes with $(u,\psi')$-adically convergent sums.

\emph{Step 1 (the contraction as a residue).} The vector $\alpha(\psi) := \eta f(-\psi)$ has its entries in $\psi^{-1}\bbQ[\psi^{-1}]\fps{u}$. As $\eta(\coef{z^{-m-1}}f(-z),e) = \coef{\psi^{-m-1}}\alpha(\psi)^Te$, the contraction is the vector $C(\psi')$ with
\[
C(\psi')^T := \Res_\psi\,\alpha(\psi)^TE_R(\psi,\psi'),
\]
the product taken in $\cA$; modulo $u^{M+1}$, $\alpha$ has finitely many non-zero coefficients, so at each order $u^M\psi'^n$ only finitely many $m$ contribute.

\emph{Step 2 (the expansion).} $E_R = (\psi+\psi')^{-1}\cN$ in $\mathfrak{gl}_5(\cA)$, so
\[
\alpha^TE_R = \alpha^T\eta^{-1}(\psi+\psi')^{-1}-\alpha^TR^{-1}(\psi)\,\eta^{-1}R^{-1}(\psi')^T(\psi+\psi')^{-1} .
\]

\emph{Step 3 (the first term).}
\[
\Res_\psi\bigl(\alpha^T\eta^{-1}(\psi+\psi')^{-1}\bigr) = \sum_j(-\psi')^j\Res_\psi\bigl(\alpha(\psi)^T\psi^{-j-1}\bigr)\eta^{-1} = \sum_j(-\psi')^j(\coef{\psi^j}\alpha)^T\eta^{-1} = 0,
\]
because $\alpha$ has no non-negative power of $\psi$.

\emph{Step 4 (the second term).} Put $\beta(\psi) := R^{-1}(\psi)^T\alpha(\psi)$, with entries in $\bbQ\lser{\psi}\fps{u}$, and $\beta_{\geq 0}(y) := \sum_{j\geq 0}y^j\coef{\psi^j}\beta$. Since $R^{-1}(\psi')^T$ does not involve $\psi$,
\[
\Res_\psi\bigl(\alpha^TR^{-1}(\psi)\,\eta^{-1}R^{-1}(\psi')^T(\psi+\psi')^{-1}\bigr) = \sum_j(-\psi')^j(\coef{\psi^j}\beta)^T\eta^{-1}R^{-1}(\psi')^T = \beta_{\geq 0}(-\psi')^T\eta^{-1}R^{-1}(\psi')^T .
\]

\emph{Step 5 (symplecticity).} The hypothesis on $R$ gives $R(\psi)^{-1} = \eta^{-1}R(-\psi)^T\eta$, so $R^{-1}(\psi)^T = \eta R(-\psi)\eta^{-1}$ ($\eta$ is symmetric) and $\beta(\psi) = \eta R(-\psi)f(-\psi) = \eta\,(Rf)(-\psi)$, where $(Rf)(-\psi)$ is the element $Rf$ of $\cH^u$ with $z$ replaced by $-\psi$. Therefore $\coef{\psi^j}\beta = (-1)^j\eta\coef{z^j}(Rf)$ and $\beta_{\geq 0}(-\psi') = \eta\sum_{j\geq 0}\psi'^j\coef{z^j}(Rf) = \eta\,(Rf)_{\geq 0}(\psi')$. With Steps 2-4,
\[
C(\psi')^T = -\bigl(\eta(Rf)_{\geq 0}(\psi')\bigr)^T\eta^{-1}R^{-1}(\psi')^T,\qquad\text{i.e.}\qquad C(\psi') = -R(\psi')^{-1}(Rf)_{\geq 0}(\psi').
\]

\emph{Step 6.} $R(\psi')^{-1}$ is invertible in $\mathfrak{gl}_5(\bbQ\fps{u,\psi'})$, so $C = 0$ if and only if $(Rf)_{\geq 0} = 0$, i.e.\ if and only if $Rf\in\cH^u_-$ (the splitting of \cref{conv:tra-forms}). The second slot follows from $E_R(\psi,\psi')^T = E_R(\psi',\psi)$.
\end{proof}

The ring matters: in $\bbQ\lser{\psi'}\lser{\psi}\fps{u}$ the split of Step 2 is a different one, and its first term has in general a non-zero residue.

\begin{cor}[Corollary E': edges in $\sKo\otimes\sKo$]\genusstatus{AI-proposed, unchecked}
\label{cor:tra-edges}
Let $P$ be a dressing with (Sp). Then (A) holds if and only if $(\Pi_0\otimes 1)E_P = 0$, and if and only if $(\Pi_i\otimes 1)E_P = (1\otimes\Pi_i)E_P = 0$ for all $i\geq 0$. If (A) holds, $E_P$ lies in $\sKo\otimes\sKo$ in the sense of \cref{lem:tra-twoslot}: $E_P = \sum_{\alpha,\beta}E_P^{\alpha\beta}\fr_\alpha(\psi)\otimes\fr_\beta(\psi')$, the sum of a summable family whose coefficients $E_P^{\alpha\beta}\in\bbQ\fps{u}$ are the flat components of $E_P$ with $a,b\neq 1$.
\end{cor}

\begin{proof}
$R^{\mathrm{vec}} = 1+O(z)$ satisfies the hypothesis of \cref{lem:tra-E}, since $R$ satisfies (Sp) (\cref{def:tra-dressing}) and (Sp) holds in vector form (\cref{conv:tra-forms}). For $k\geq 1$, $z^{-k}J^{\mathrm C}\in\cH^u_-$ (\cref{lem:tra-conv}~(c)). In vector form $P^{\mathrm{vec}}J^{\mathrm C} = \eta^{-1}P\eta\,\eta^{-1}\YC = \eta^{-1}(P\YC) = \eta^{-1}Y'$; as $P_0$ does not involve $z$ and multiplication by $z^{-k}$ shifts the splitting,
\[
(R^{\mathrm{vec}}z^{-k}J^{\mathrm C})_{\geq 0} = z^{-k}\,(P_0^{\mathrm{vec}})^{-1}\eta^{-1}\,(Y')_{\geq k} .
\]
If (A) holds, this vanishes for every $k\geq 1$, and \cref{lem:tra-E} shows that $E_P$ annihilates $\Omega_{\mathrm s}(z^{-k}J^{\mathrm C},\cdot) = (-1)^{k+1}\Pi_{k-1}$ \eqref{eq:tra-pairing} in each slot, for every $k\geq 1$; this is $(\Pi_i\otimes 1)E_P = (1\otimes\Pi_i)E_P = 0$ for all $i\geq 0$. Conversely, if $(\Pi_0\otimes 1)E_P = 0$, then \cref{lem:tra-E} with $f = z^{-1}J^{\mathrm C}$ gives $(Y')_{\geq 1} = 0$, which is (A). The frame expansion is \cref{lem:tra-twoslot}.
\end{proof}

\begin{lem}[Lemma T: the leaf point]\genusstatus{AI-proposed, unchecked}
\label{lem:tra-T}
Let $P$ be a dressing with (Sp) and (A).
\begin{description}
\item[(a)] $\Omega_{\mathrm s}(f,q_T) = \Omega_{\mathrm s}(P^{\mathrm{vec}}f,-z1)$ for every $f\in\cH^u$.
\item[(b)] $\Pi_0(q_T) = T_P$ and $\Pi_i(q_T) = 0$ for $i\geq 1$; so $q_T\in\sK$.
\item[(c)] With $T_{\geq 2}\in z^2\cH^u_+$ as in \eqref{eq:tra-leaf}, $q_T = -(z/d)1+T_{\geq 2}$, and $\tilde q_T = d\,q_T = -z1+dT_{\geq 2}$ lies in $-z+\sK$, with
\[
\Pi_0(\tilde q_T) = d\,T_P =: v+\delta_P .
\]
Moreover $dT_{\geq 2} = \tilde q_T+z1 = \ttr_{R^{\mathrm{vec}}}(z)$ lies in $\sK\cap z^2\cH^u_+$, and $dT_{\geq 2} = y^\circ+\delta_PH$ with $y^\circ\in\sKo$, whose frame coordinates $y^\circ_{(k,a)}$ ($a\neq 1$) vanish for $k\leq 1$.
\item[(d)] $\delta_P\in u^2\bbQ\fps{u}$; more precisely $T_P-p_{0,0}v\in u^2\bbQ\fps{u}$ and $d\,p_{0,0} = 1$.
\end{description}
\end{lem}

\begin{proof}
(a) By (Sp) in vector form, $\eta\,(P^{\mathrm{vec}})^{-1}(z) = P^{\mathrm{vec}}(-z)^T\eta$; hence $\eta(f(-z),(P^{\mathrm{vec}})^{-1}(z)g(z)) = \eta((P^{\mathrm{vec}}f)(-z),g(z))$ in $\bbQ\lser{z}\fps{u}$, and $\Omega_{\mathrm s}(f,(P^{\mathrm{vec}})^{-1}g) = \Omega_{\mathrm s}(P^{\mathrm{vec}}f,g)$ for $f,g\in\cH^u$, by associativity in $\mathfrak{gl}_5(\bbQ\lser{z}\fps{u})$ (\cref{lem:tra-conv}~(c)). Take $g = -z1$.

(b) By \eqref{eq:tra-pairing} with $k = i+1$, by (a) and by $P^{\mathrm{vec}}J^{\mathrm C} = \eta^{-1}Y'$,
\[
\Pi_i(q_T) = (-1)^i\,\Omega_{\mathrm s}(z^{-i-1}J^{\mathrm C},q_T) = (-1)^i\,\Omega_{\mathrm s}(z^{-i-1}\eta^{-1}Y',-z1).
\]
Since $(\eta^{-1}Y')^T\eta\,1 = Y'_0$, the last pairing is $\Res_z(-z)^{-i}Y'_0(-z)\,\dd z = -\coef{z^{i-1}}Y'_0$; so
\[
\Pi_i(q_T) = (-1)^{i+1}\coef{z^{i-1}}Y'_0 .
\]
For $i = 0$ this is $T_P$; for $i\geq 2$ it vanishes by (A). For $i = 1$ it vanishes by the grading: $(P\YC)_0 = \sum_kz^kp_{k,-k}\YC_{-k}$, and $\YC_{-k}$ has only the powers $z^{-n}$ with $n\equiv\nu_{-k}+1\equiv 1+k\pmod 5$ (\cref{lem:gau-vanshape}), so $Y'_0$ has only the powers $z^{-n}$ with $n\equiv 1\pmod 5$, and $\coef{z^0}Y'_0 = 0$.

(c) The formulas for $q_T$, $\tilde q_T$ and $dT_{\geq 2}$ are \eqref{eq:tra-leaf}. Since $\sK$ is a $\bbQ\fps{u}$-module and $-z+\sK = \sK$ (\cref{cor:tra-slice}), (b) gives $\tilde q_T\in -z+\sK$ with $\Pi_0(\tilde q_T) = dT_P$. As $-z1\in\sK$ (\cref{cor:tra-slice}), $dT_{\geq 2} = \tilde q_T+z1$ lies in $\sK\cap z^2\cH^u_+$, with $\Pi_0(dT_{\geq 2}) = dT_P+\Pi_0(z1) = dT_P-v = \delta_P$, and \cref{cor:tra-slice} with \cref{lem:tra-conv}~(b3) gives the decomposition: $y^\circ = dT_{\geq 2}-\delta_PH$ has the same flat components with $a\neq 1$ as $dT_{\geq 2}\in z^2\cH^u_+$.

(d) By the grading, $T_P = -\coef{z^{-1}}(P\YC)_0 = -\sum_{k\geq 0}p_{k,-k}\coef{z^{-1-k}}\YC_{-k}$. The term $k = 0$ is $-p_{0,0}\coef{z^{-1}}\YC_0 = p_{0,0}v$, as $\coef{z^{-1}}\YC_0 = -v$ \eqref{eq:tra-values}. For $k\geq 1$ the coefficient $\coef{z^{-1-k}}\YC_{-k}$ has $u$-order at least $(k+1)-\lfloor(k+1)/5\rfloor\geq 2$ (\cref{lem:tra-conv}~(c)), and $p_{k,-k}\in\bbQ\fps{u}$; so $T_P-p_{0,0}v\in u^2\bbQ\fps{u}$. By (Sp) at $z^0$, $p_{0,0}\,\eta_{03}\,p_{0,3} = \eta_{03}$, i.e.\ $d\,p_{0,0} = 1$. Hence $dT_P = v+d\,(T_P-p_{0,0}v)$ and $\delta_P = dT_P-v\in u^2\bbQ\fps{u}$.
\end{proof}

\begin{lem}[Lemma F: finiteness]\genusstatus{AI-proposed, unchecked}
\label{lem:tra-F}
Let $P$ be a dressing with (Sp) and $\Gamma$ a leg-free stable graph. The contraction along $\Gamma$ of the edges $E_P$ with the vertex tensors $d^{2g_v-2+n_v}\dd^{n_v}\overline F_{g_v}(\tilde q_T)$, one slot per half-edge (this is $|\Aut\Gamma|\,\mathrm{Cont}_\Gamma$ in the notation of \cref{def:red-flatsum}), is a finite sum of products of flat model correlators at $v = v(u)$, which are monomials $Cv(u)^e$ with $C\in\bbQ$ and $e\in\bbZ$, and of polynomials in the entries of the $P_k$, the $(P^{-1})_k$ and $d^{\pm 1}$. If $d(0) = 1$, each vertex tensor equals the Taylor expansion $\dd^{n_v}\overline F_{g_v}(q_T)$ of $\overline F_{g_v}$ at the base point $-z1$ in the direction $q_T+z1$.
\end{lem}

\begin{proof}
\emph{Step 1 (vertices).} At a vertex of genus $h$ with $n$ half-edges, $2h-2+n>0$, the definition of $\overline F_h$ (\cref{def:set-graphsum,lem:set-dilaton}) and $\tilde q_T+z1 = dT_{\geq 2}$ \eqref{eq:tra-leaf} give
\[
d^{2h-2+n}\,\dd^n\overline F_h(\tilde q_T)[X_1,\dots,X_n] = d^{2h-2+n}\sum_{m\geq 0}\frac1{m!}\langle X_1\cdots X_n(dT_{\geq 2})^m\rangle_h .
\]
If $d(0) = 1$, \cref{lem:tra-conv}~(f1) with $\varkappa = d$ and $y = T_{\geq 2}$ shows that this equals the $u$-adically convergent Taylor expansion $\sum_m\frac1{m!}\langle X_1\cdots X_n(q_T+z1)^m\rangle_h = \dd^n\overline F_h(q_T)[X_1,\dots,X_n]$; at $h = 1$ this needs $n\geq 1$, which holds at every genus-$1$ vertex of a leg-free stable graph of genus at least $2$, so the logarithmic genus-$1$ term of the dilaton equation does not enter. \emph{Step 2 (leaves).} Since $dT_{\geq 2}\in z^2\cH^u_+$, the sum over $m$ stops at $m = 3h-3+n$, and only slot degrees $\leq 3h-3+n$ occur (\cref{lem:tra-conv}~(d)). \emph{Step 3 (edges).} So only finitely many coefficients $\coef{\psi^k\psi'^l}E_P = \sum_{j=0}^l(-1)^j\coef{\psi^{k+1+j}\psi'^{l-j}}\cN$ of each edge matter, and $\cN$ is a polynomial in the entries of $(R^{\mathrm{vec}})^{-1} = (P^{\mathrm{vec}})^{-1}P_0^{\mathrm{vec}}$; the coefficients of $dT_{\geq 2} = z1-zd\,(P^{\mathrm{vec}})^{-1}(z)1$ \eqref{eq:tra-leaf} are polynomials in $d$ and in the entries of the $(P^{-1})_k$. The flat model correlators at $v(u)$ are monomials $Cv(u)^e$ by \cref{lem:set-modhom}.
\end{proof}

\begin{prop}[the canonical gauge as a dressing]\genusstatus{AI-proposed, unchecked}
\label{prop:tra-can}
Let $(\alpha,\beta)\in\bbQ^2$ give a conifold frame (\cref{def:set-frames}) and let $v(u) = u\tau(u)^{-5/4}$ (\cref{def:set-tau}), so that $v(0) = 0$ and $v'(0) = 1$. The canonical gauge $\Pc$ (\cref{def:gau-can}) is a dressing of the model at $v(u)$ with (Sp) and (A), with $Y' = \YCQ$, $T_{\Pc} = \tc$ and $d(0) = 1$. Hence the edges of $\Pc$ lie in $\sKo\otimes\sKo$, as $\bbQ\fps{u}$-combinations of frames (\cref{cor:tra-edges}); $q_T\in\sK$ with $\Pi_0(q_T) = \tc$, and $\tilde q_T = d\,q_T\in -z+\sK$ with $\Pi_0(\tilde q_T) = d\,\tc = v+\delta_{\Pc}$, $\delta_{\Pc}\in u^2\bbQ\fps{u}$ (\cref{lem:tra-T}); and every graph term is a finite sum, each vertex tensor being the Taylor expansion at the base point (\cref{lem:tra-F}).
\end{prop}

\begin{proof}
By \cref{thm:gau-U}~(d) and \cref{def:gau-can}, $\Pc = \sum_kz^kP_k$ has $P_k\in\mathfrak{gl}_5(\bbQ\fps{u})$, is graded and has $P_0$ invertible: it is a dressing of the model at $v(u)$, the model being pulled back by $v(u)$ as in \cref{prop:gau-equalprod}. (Sp) is \cref{prop:gau-symp}. By \cref{thm:gau-G}~(a), $\Pc\YC(v(u)) = \YCQ$, which lies in $\bbQ[z^{-1}]\fps{u}^5$ by \cref{lem:gau-vanq}~(c), so it has no positive power of $z$: this is (A), with $Y' = \YCQ$. By \cref{lem:gau-vanq}~(d), $\coef{z^{-1}}\YCQ = -\tc\,\mathbf e_0$, so $T_{\Pc} = \tc$. By \cref{lem:gau-diag}, which applies because (N0) holds in the conifold frames (\cref{def:set-frames,prop:gau-equalprod}), $p_{0,a}(0) = 1$ for all $a$; so $d = p_{0,3} = 1/p_{0,0}$ (\cref{lem:tra-T}~(d)) has $d(0) = 1$. The rest is \cref{cor:tra-edges,lem:tra-twoslot,lem:tra-T,lem:tra-F}.
\end{proof}



